\usepackage[margin=1in]{geometry}
\usepackage{amsmath,amssymb,amsthm,mathtools}
\usepackage{enumitem}
\usepackage{hyperref}
\usepackage{tikz}
\usetikzlibrary{arrows.meta,calc,positioning,decorations.markings,decorations.pathmorphing,
  patterns,shadings,shapes.geometric,fadings,backgrounds,intersections,bending}

\definecolor{formteal}{HTML}{1f4e79}
\definecolor{formgold}{HTML}{b88a14}
\definecolor{formrust}{HTML}{a8431d}
\definecolor{formplum}{HTML}{6b2e6e}
\definecolor{formsage}{HTML}{467b3a}
\definecolor{formslate}{HTML}{3a5168}
\definecolor{formivory}{HTML}{f6efe0}

% Location of the Sage-generated coordinate files, relative to the directory
% pdflatex is run from. Compile the book from inside book/:
%     cd book && pdflatex main.tex
% If you compile from elsewhere, redefine \sagedir to point at ../sagemath.
\newcommand{\sagedir}{../sagemath}

% Sage-computed precise coordinates for every TikZ figure.
% Generated by sagemath/build_coords.sh.
\input{\sagedir/coords.tex}

% Precise torus rendering computed by sagemath/08_torus_render.sage.
% Includes the outer and inner silhouettes plus six meridians and four
% latitude circles, each split into visible/hidden arcs.
\input{\sagedir/coords_torus.tex}

% Precise elliptic curve y^2 = 4x^3 - 4x sample points.
% Computed by sagemath/12_curve_sample.sage.
% Provides \curveOvalFull (smooth-cycle), \curveOvalUpper, \curveOvalLower,
% \curveBranchUpper, \curveBranchLower (all clipped at y = \curveYclip).
\input{\sagedir/coords_curve_sample.tex}

% Convenience macro that paints the precise torus inside the current scope.
% Usage:  \begin{scope}[shift={(...)}, scale=...] \drawprecisetorus \end{scope}
% The sage script computes for tilt angle beta = 60 deg:
%   - outer silhouette  (\torusOuter, closed curve)
%   - inner silhouette  (\torusInner, closed curve = donut-hole outline)
%   - 8 meridians at v = k pi/4 (small circles), split into visible/hidden arcs
%   - 4 longitudes at u in {0, pi/2, pi, 3pi/2} (latitude circles),
%     split into visible/hidden arcs
\newcommand{\drawprecisetorus}{%
  % (1) Fill the visible torus body with light teal.
  \fill[formteal!13, line join=round]
    plot[smooth cycle] coordinates {\torusOuter};
  % (2) Cut out the donut hole using the inner silhouette curve.
  \fill[white, line join=round]
    plot[smooth cycle] coordinates {\torusInner};
  % (3) Visible meridian arcs (eight small circles spaced 45 deg apart).
  % Back meridians (v in [0, pi]) form one visible arc each;
  % front meridians (v in (pi, 2 pi)) wrap across u = 0 and are split
  % into two visible chunks (VA and VB) by the sampling boundary.
  \draw[draw=formteal!60!black, line width=0.42pt, line join=round, opacity=0.9]
    plot[smooth] coordinates {\torusMerZeroVA};
  \draw[draw=formteal!60!black, line width=0.42pt, line join=round, opacity=0.9]
    plot[smooth] coordinates {\torusMerOneVA};
  \draw[draw=formteal!60!black, line width=0.42pt, line join=round, opacity=0.9]
    plot[smooth] coordinates {\torusMerTwoVA};
  \draw[draw=formteal!60!black, line width=0.42pt, line join=round, opacity=0.9]
    plot[smooth] coordinates {\torusMerThreeVA};
  \draw[draw=formteal!60!black, line width=0.42pt, line join=round, opacity=0.9]
    plot[smooth] coordinates {\torusMerFourVA};
  \draw[draw=formteal!60!black, line width=0.42pt, line join=round, opacity=0.9]
    plot[smooth] coordinates {\torusMerFiveVA};
  \draw[draw=formteal!60!black, line width=0.42pt, line join=round, opacity=0.9]
    plot[smooth] coordinates {\torusMerFiveVB};
  \draw[draw=formteal!60!black, line width=0.42pt, line join=round, opacity=0.9]
    plot[smooth] coordinates {\torusMerSixVA};
  \draw[draw=formteal!60!black, line width=0.42pt, line join=round, opacity=0.9]
    plot[smooth] coordinates {\torusMerSixVB};
  \draw[draw=formteal!60!black, line width=0.42pt, line join=round, opacity=0.9]
    plot[smooth] coordinates {\torusMerSevenVA};
  \draw[draw=formteal!60!black, line width=0.42pt, line join=round, opacity=0.9]
    plot[smooth] coordinates {\torusMerSevenVB};
  % (4) Visible longitude arcs (the four canonical latitude circles).
  % Outer (u=0) is visible only on the front half (v in (pi, 2 pi)),
  % inner (u=pi) only on the back half (v in (0, pi)); the top crest
  % (u=pi/2) is visible everywhere; the bottom (u=3pi/2) is fully hidden.
  \draw[draw=formteal!60!black, line width=0.5pt, line join=round, opacity=0.95]
    plot[smooth] coordinates {\torusLonTopVA};
  \draw[draw=formteal!60!black, line width=0.42pt, line join=round, opacity=0.8]
    plot[smooth] coordinates {\torusLonOuterVA};
  \draw[draw=formteal!60!black, line width=0.42pt, line join=round, opacity=0.8]
    plot[smooth] coordinates {\torusLonInnerVA};
  % (5) Crisp silhouette outlines drawn last so they sit on top.
  \draw[draw=formteal!65!black, line width=0.95pt, line join=round]
    plot[smooth cycle] coordinates {\torusOuter};
  \draw[draw=formteal!65!black, line width=0.8pt, line join=round]
    plot[smooth cycle] coordinates {\torusInner};
}

% Precise sphere (CP^1) rendering, matching the torus style.
% Computed by sagemath/10_sphere_render.sage at gamma = 30 deg elevation.
% Provides:
%   \sphereOutline    -- silhouette as a closed coordinate list (unit circle)
%   \sphereLatXY{V,H}A -- equator/tropics/etc. visible/hidden arcs
%   \sphereMerXyz{V,H}A -- eight meridians at v = k*pi/4
%   \sphereNorth{x,y}, \sphereSouth{x,y} -- pole projection coordinates
\newcommand{\drawprecisesphere}{%
  % Color tone matches the torus (formteal) for visual unity across figures.
  % (1) Flat-filled body of the sphere (same filling style as the torus)
  \fill[formteal!13, line join=round]
    plot[smooth cycle] coordinates {\sphereOutline};
  % (2) Visible latitude arcs
  \draw[draw=formteal!60!black, line width=0.45pt, line join=round, opacity=0.9]
    plot[smooth] coordinates {\sphereLatEquatorVA};
  \draw[draw=formteal!60!black, line width=0.4pt, line join=round, opacity=0.85]
    plot[smooth] coordinates {\sphereLatNorthVA};
  \draw[draw=formteal!60!black, line width=0.4pt, line join=round, opacity=0.85]
    plot[smooth] coordinates {\sphereLatNorthVB};
  \draw[draw=formteal!60!black, line width=0.4pt, line join=round, opacity=0.85]
    plot[smooth] coordinates {\sphereLatTropicVA};
  \draw[draw=formteal!60!black, line width=0.4pt, line join=round, opacity=0.85]
    plot[smooth] coordinates {\sphereLatTropicVB};
  \draw[draw=formteal!60!black, line width=0.4pt, line join=round, opacity=0.85]
    plot[smooth] coordinates {\sphereLatSouthVA};
  % (3) Visible meridian arcs (eight meridians at v = k pi/4)
  \draw[draw=formteal!60!black, line width=0.42pt, line join=round, opacity=0.9]
    plot[smooth] coordinates {\sphereMerZeroVA};
  \draw[draw=formteal!60!black, line width=0.42pt, line join=round, opacity=0.9]
    plot[smooth] coordinates {\sphereMerOneVA};
  \draw[draw=formteal!60!black, line width=0.42pt, line join=round, opacity=0.9]
    plot[smooth] coordinates {\sphereMerTwoVA};
  \draw[draw=formteal!60!black, line width=0.42pt, line join=round, opacity=0.9]
    plot[smooth] coordinates {\sphereMerThreeVA};
  \draw[draw=formteal!60!black, line width=0.42pt, line join=round, opacity=0.9]
    plot[smooth] coordinates {\sphereMerFourVA};
  \draw[draw=formteal!60!black, line width=0.42pt, line join=round, opacity=0.9]
    plot[smooth] coordinates {\sphereMerFiveVA};
  \draw[draw=formteal!60!black, line width=0.42pt, line join=round, opacity=0.9]
    plot[smooth] coordinates {\sphereMerSixVA};
  \draw[draw=formteal!60!black, line width=0.42pt, line join=round, opacity=0.9]
    plot[smooth] coordinates {\sphereMerSevenVA};
  % (4) Hidden back arcs as dashed light strokes
  \draw[draw=formteal!50!black, dashed, line width=0.3pt, line join=round, opacity=0.55]
    plot[smooth] coordinates {\sphereLatEquatorHA};
  % (5) Crisp silhouette outline on top
  \draw[draw=formteal!65!black, line width=0.95pt, line join=round]
    plot[smooth cycle] coordinates {\sphereOutline};
}

% Convenience macro that draws the real locus of E: y^2 = 4x^3 - 4x
% using pre-computed coordinate lists from sagemath/12_curve_sample.sage.
% Usage:  \begin{scope}[shift={(...)}, scale=...] \drawprecisecurve \end{scope}
% The scope should set up an appropriate coordinate system.
% This macro draws only the curve arcs and branch-point dots; the caller
% is responsible for axes and labels.
\newcommand{\drawprecisecurve}{%
  % Oval (compact component, x in [-1,0]): draw as closed smooth cycle.
  \draw[draw=formteal!60!black, line width=0.88pt, line join=round]
    plot[smooth cycle, tension=0.85] coordinates {\curveOvalFull};
  % Branch (unbounded component, x >= 1): upper and lower arcs.
  \draw[draw=formteal!60!black, line width=0.88pt, line join=round]
    plot[smooth, tension=0.9] coordinates {\curveBranchUpper};
  \draw[draw=formteal!60!black, line width=0.88pt, line join=round]
    plot[smooth, tension=0.9] coordinates {\curveBranchLower};
  % Branch points (Weierstrass roots e1=-1, e2=0, e3=1)
  \fill[formteal!70!black] (-1,0) circle (1.8pt);
  \fill[formteal!70!black] ( 0,0) circle (1.8pt);
  \fill[formteal!70!black] ( 1,0) circle (1.8pt);
}

\setlist[enumerate]{leftmargin=*, itemsep=0.35em}
\setlist[itemize]{leftmargin=*, itemsep=0.25em}
\allowdisplaybreaks
\emergencystretch=2em

\newcommand{\C}{\mathbb{C}}
\newcommand{\R}{\mathbb{R}}
\newcommand{\Z}{\mathbb{Z}}
\newcommand{\Q}{\mathbb{Q}}
\newcommand{\F}{\mathbb{F}}
\newcommand{\CP}{\mathbb{CP}}
\newcommand{\HH}{\mathbb{H}}
\newcommand{\OO}{\mathcal{O}}
\newcommand{\M}{\mathcal{M}}
\newcommand{\Div}{\operatorname{Div}}
\newcommand{\Pic}{\operatorname{Pic}}
\newcommand{\Prin}{\operatorname{Prin}}
\newcommand{\Jac}{\operatorname{Jac}}
\newcommand{\smap}{\mathfrak{s}}% Abel--Jacobi sum map Div^0(T) -> T (distinct from Weierstrass sigma)
\newcommand{\Span}{\operatorname{span}}
\newcommand{\divisor}{\operatorname{div}}
\newcommand{\ord}{\operatorname{ord}}
\newcommand{\id}{\operatorname{id}}
\newcommand{\Res}{\operatorname{Res}}
\newcommand{\degdiv}{\operatorname{deg}}
\newcommand{\dd}{\,d}
\newcommand{\wpprime}{\wp'}
\newcommand{\LambdaLat}{\Lambda}

% Textbook theorem environments: all share a single counter that resets each
% chapter (Theorem 2.1, Proposition 2.2, Definition 2.3, ...) so statements read
% in one sequence. thmtools' `sibling=theorem' gives the shared numbering, while
% each environment keeps its OWN label type, so cleveref prints the right word
% (\cref to a Proposition reads "proposition 2.2", not "theorem 2.2").
% Environment NAMES match the source notes, so reused content drops in unchanged.
\usepackage{thmtools}
\declaretheorem[name=Theorem,numberwithin=chapter]{theorem}
\declaretheorem[name=Proposition,sibling=theorem]{proposition}
\declaretheorem[name=Lemma,sibling=theorem]{lemma}
\declaretheorem[name=Corollary,sibling=theorem]{corollary}
\declaretheorem[name=Definition,sibling=theorem,style=definition]{definition}
\declaretheorem[name=Example,sibling=theorem,style=definition]{example}
\declaretheorem[name=Remark,sibling=theorem,style=remark]{remark}
% `warning' cannot use that environment name directly inside thmtools, because
% the internal \@warning is already reserved (by amsmath/mathtools). We declare
% it under the safe name `warningthm' (sibling => shared theorem counter, so it
% joins the Theorem/Proposition/... sequence) and alias \warning to it, so
% authors still write \begin{warning} ... \end{warning}.
\declaretheorem[name=Warning,sibling=theorem,style=remark]{warningthm}
\let\warning\warningthm
\let\endwarning\endwarningthm
% Exercises run on their own per-chapter counter.
\declaretheorem[name=Exercise,numberwithin=chapter,style=definition]{exercise}

\newenvironment{keyidea}[1]{%
  \par\medskip
  \begin{center}
  \begin{minipage}{0.92\textwidth}
  \hrule
  \smallskip
  \textbf{#1.}\quad\ignorespaces
}{%
  \smallskip
  \hrule
  \end{minipage}
  \end{center}
  \medskip
}

\tikzset{
  arrow/.style={-{Latex[length=2mm]}, thick},
  point/.style={circle, fill=black, inner sep=1.5pt},
  labelnode/.style={font=\small, align=center},
  note/.style={draw=black!55, rounded corners=2pt, align=center, inner sep=5pt, fill=black!3},
  visualnode/.style={draw=black!60, rounded corners=2pt, align=center, inner sep=6pt, fill=black!2}
}

% Smart cross-references (\cref / \Cref). Load last, after hyperref.
\usepackage{cleveref}
% cleveref has built-in names for theorem/proposition/lemma/corollary/
% definition/remark/example, so those resolve automatically. The two custom
% types -- `warning' and `exercise' -- have no built-in, so name them here.
% Lowercase singular is for \cref ("exercise 2.3"); capitalized for \Cref.
\crefname{warningthm}{warning}{warnings}
\Crefname{warningthm}{Warning}{Warnings}
\crefname{exercise}{exercise}{exercises}
\Crefname{exercise}{Exercise}{Exercises}
